%!TEX root = ../template.tex
%%%%%%%%%%%%%%%%%%%%%%%%%%%%%%%%%%%%%%%%%%%%%%%%%%%%%%%%%%%%%%%%%%%
%% chapter1.tex
%% NOVA thesis document file
%%
%% Chapter with introduction
%%%%%%%%%%%%%%%%%%%%%%%%%%%%%%%%%%%%%%%%%%%%%%%%%%%%%%%%%%%%%%%%%%%

\typeout{NT FILE chapter1.tex}%

\chapter{Introdução}
\label{cha:introducao}

Com o crescimento do controlo político dos media, o uso de ferramentas anónimas para contacto com a world wide web torna-se um método viável para obter notícias sobre o mundo sem restrições. Uma das principais ferramentas usadas é o Tor, desenvolvido pelo Tor Project, que permite a navegação na web de forma anónima, protegendo a privacidade e ocultando a identidade do utilizador \cite{tor}. Como explorado por Peter Story et al, quando uma amostra de utilizadores do Tor são questionados sobre o seu uso, os resultados ditam que a procura de informação sobre políticas, celebridades, eventos atuais, ficheiros partilhados por lapso, entre outros, está nos primeiros cinco motivos mais comuns para o uso da ferramenta \cite{story2022increasing}. Há cerca de oito anos a Turquia foi, por exemplo, identificada como o país com mais pedidos de censura de informação na rede do Twitter (agora X), emitidos em grande parte pelas autoridades turcas. No seguimento de ataques terroristas em Suruç, Ankara e Istambul, o governo turco fez pedidos de remoção de imagens de vítimas destes ataques, que com análise \textit{a posteriori} envolveu uma censura muito maior do que o esperado, removendo também conteúdo que criticava o governo turco \cite{sandberg2017geography}. O uso do Tor como ferramenta para obter notícias sem censura é usado como heurística para detetar o nível de censura do país em questão, onde um uso mais elevado geralmente leva a um maior nível de censura \cite{sandberg2017geography}.

Ao longo das últimas duas décadas, surgiram novas métricas e métodos de prevenção que visam supervisionar a rede Tor para eliminar possíveis redes de \textit{sybils} \cite{levine2006survey}. Este trabalho visa construir uma nova forma de prevenção deste tipo de ataques, alterando o paradigma de “procurar e analisar possíveis ataques” para “impossibilitar a exploração da vulnerabilidade”, por meio de ferramentas de identificação anónima. Este tipo de ferramentas foi explorado em 2025 a partir de \textit{Zero-Knowledge Proofs}, que permitem a verificação de identidades e atributos sem exposição de informação sensível, combinando-as com o uso de DIDs, \textit{Decentralized Identifiers}, para formar \textit{frameworks} para identidade descentralizada sem revelar o utilizador \cite{yuan2025scalable}, evoluindo para protocolos de identificação descentralizados, também em 2025 \cite{liu2025diap}. O objetivo deste trabalho é o desenvolvimento de um esquema que permita identificação descentralizada de utilizadores na rede Tor, garantindo proteção contra \textit{Sybil attacks} \cite{douceur2002sybil} numa modalidade que não revele informação relativa a nenhuma das entidades e proteja a autonomia sobre a identidade do cliente e os atributos que pretende divulgar.

\section*{Requisitos da Solução}
\label{sec:requisitos_solucao}

A identificação anónima de um novo nó dentro da rede Tor surge como tentativa de prevenção de possíveis ataques \textit{Sybil} \cite{douceur2002sybil}, ao autenticar a sua identidade às \textit{Directory Authorities} do Tor sem revelar nenhuma informação que o identifique. Para tal, é necessário uma autenticação que o categoriza como confiável, que deve ser definida para efeitos de avaliação da solução. Para além disto, a solução deve proteger a identidade do dono do nó, confirmando-se deste modo anónima e inútil para efeitos de uso para derivar novas informações relativas ao mesmo ou, por outras palavras, \textit{zero-knowledge}. Por último, é também necessário que uma autenticação feita por um \textit{sybil} seja sempre detetada por meio de autenticação (DA), de forma a que não exista forma de contornar a verificação.

Para que a solução proposta seja avaliada de acordo com a sua validade, é necessário a criação dos seguintes requisitos de avaliação:

\begin{itemize}
  \item \textbf{\textit{Zero-Knowledge:}} Os donos dos nós devem poder conseguir provar a sua credibilidade para correr um nó na rede Tor sem ser necessário revelar a sua identidade, nem derivar informações que possam levar ao mesmo efeito;
  \item \textbf{\textit{Confiabilidade:}} Um novo utilizador é considerado confiável se apresentar uma prova a partir de terceiros confiáveis que valide a sua identidade ou se tiver algo em risco de perda caso seja classificado como malicioso; 
  \item \textbf{\textit{Obrigatoriedade de Autenticação:}} Não deve ser possível, para um novo utilizador na rede, evadir a autenticação feita à sua identidade, nem apresentar um passaporte digital ou certificado que induza em erro à sua confiabilidade;
  \item \textbf{\textit{Repelente de Sybils:}} A solução deve mitigar \textit{Sybil attacks}  (mais sobre este tipo de ataque no capítulo do Contexto), de forma a que deixe de existir uma ameaça à rede Tor a partir de múltiplas identidades para uma só entidade;
  \item \textbf{\textit{Requisito de Performance:}} A performance da implementação deve ser adequada para o uso corrente no dia-a-dia da rede Tor, atingindo os mínimos requisitos necessários e contemplando todos os mecanismos possíveis para melhorá-la.
\end{itemize}

A partir destas definições, é possível aplicar um método científico de avaliação da solução proposta, podendo concluir resultados fidedignos e absolutos. 

Esta proposta da tese está estruturada de forma a que o leitor, mesmo que leigo na matéria, tenha o conhecimento mínimo necessário no decorrer da leitura para entender todos os conceitos fulcrais da solução proposta, partindo da explicação do Contexto e Trabalho Relacionado para a solução e a metodologia proposta, terminando na avaliação da mesma de acordo com os requisitos enunciados anteriormente.